\subsection{A uniform Cartier calculation for every even product type}
\label{sec:DE-A3-all-ss-rectangles}

\begin{lemma}[A rank-two Frobenius bundle and its endomorphism quotient]
\label{lem:DE-A3-ss-Frobenius-endomorphisms}
\status{Proved; root, reduction-descent, and monodromy full independent proof and certificate audits passed.}
Let $E/\mathbf F_2$ be $y^2+y=x^3$, let $O$ be its point at infinity,
and identify its Frobenius twist with $E$ by this model.  Write $F$ for
relative Frobenius and $\partial x=1$, $\partial y=x^2$ for the invariant
vector field.  For $q>0$ put
\[
 V=F_*\OO_E(2qO),\qquad H=\OO_E(qO).
\]
Then $V\simeq U\otimes H$, where $U=F_*\OO_E$ is the nonsplit extension
of $\OO_E$ by itself.  There is a perfect alternating pairing
$V\otimes V\to\det V\simeq\OO_E(2qO)$, given by Cartier trace with
respect to $dx$.  Moreover
\[
 h^0(E,\End V)=2,\qquad
 h^0\bigl(E,\End V/\OO_E\operatorname{id}\bigr)=2.
\]
The latter two-dimensional space is spanned by the classes of
multiplication by the source function $x$ and of $\partial$.
\end{lemma}
\begin{proof}
The Cartier formula $C(dx)=0$ makes Frobenius zero on $H^1(E,\OO_E)$
by Serre duality.  In
$0\to\OO_E\to F_*\OO_E\to B^1_E\to0$, the quotient is a line
bundle of degree zero with one nonzero global section.  It is therefore
$\OO_E$.  The extension is nonsplit because $h^0(F_*\OO_E)=1$.
Projection formula gives $V=U\otimes H$ and $\det V=H^2$.
We use Cartier in the convention of
\cite[0, Theorem~2.1.9, p.~515]{Illusie1979DRW}, already opened and
checked in the crystalline comparison above.
The pairing in the rational basis $(1,x)$ is the wedge pairing:
$C(dx)=C(x^2dx)=0$, while $C(xdx)=dx$.  It is thus alternating and
perfect, and the Cartier transformation law makes it global.

The dual of the displayed nonsplit extension is again nonsplit, so
$\dim\Hom(U,\OO_E)=1$.  Applying $\Hom(U,-)$ to the extension,
the identity maps to its nonzero quotient homomorphism.  Consequently
$\dim\End(U)=2$.  The identity and the nonzero nilpotent endomorphism
$\partial$ span this space; $\partial^2=0$.  Both have trace zero
in characteristic two.  The trace pairing makes $\End(U)$ self-dual.
Thus the map $H^1(\OO_E)\to H^1(\End(U))$ induced by the scalar
inclusion is dual to the zero map on global traces.  The cohomology
sequence of $0\to\OO_E\to\End(U)\to\End(U)/\OO_E\to0$ now gives
the second dimension assertion.

The derivative is a global endomorphism of $V$: the line bundle
$\OO_E(2qO)=F^*H$ has horizontal square transition functions.
Multiplication by $x$ gives a global class modulo scalar endomorphisms,
although multiplication by $x$ itself has a pole at $O$.  Indeed, for
the uniformizer $t=x/y$ at $O$ the equation of $E$ gives the exact identity
\[
                    x+t^{-2}=y/x^2.
\]
The right side is regular at $O$, and multiplication by $t^{-2}$ is a
scalar endomorphism over the Frobenius base because it is a square.
Away from $O$, multiplication by $x$ is regular.  In the rational basis
$(1,x)$, with $X$ denoting the coordinate on the Frobenius base, the two
classes have representatives
\[
 \begin{pmatrix}0&X\\1&0\end{pmatrix},\qquad
 \begin{pmatrix}0&1\\0&0\end{pmatrix}.
\]
They are linearly independent modulo scalars.  The preceding dimension
calculation shows that they span all global classes.  This argument
uses no classification of vector bundles on elliptic curves.
\end{proof}

\begin{lemma}[A uniform sparse section]
\label{lem:DE-A3-ss-uniform-sparse}
\status{Proved; root, reduction-descent, and monodromy full independent proof and certificate audits passed.}
For every positive integer $q$ divisible by six, put
$N_q=\OO_E(2qO)\boxtimes\OO_E(3O)$ on $E\times E$.
In coordinates $(x,y)$ and $(z,w)$ take the section
\[
\begin{split}
 s_q&=A+Bz+Cw,\\
 A&=x+P\{x^2+(1+x)y\},\\
 B&=x+Py^2,\\
 C&=1+Py,\\
 P&=x^{q-3}.
\end{split}
\]
All coefficients lie in $\mathbf F_2$.  The root-removed Cartier
endomorphism $C_{s_q}$ of $H^0(E\times E,N_q)$ satisfies
\[
             \rank C_{s_q}=\rank C_{s_q}^2=6q-2.
\]
In particular a general smooth member of $|N_q|$ has full and relative
stable rank $6q-2$.
\end{lemma}
\begin{proof}
The pole orders of $A,B,C$ at $O$ are respectively $2q-1$, $2q$, and
$2q-3$, so the displayed functions are sections as asserted.  Set
$u=x^{q/2-2}$; then $P=xu^2$.  Reduction by $y^2+y=x^3$ gives
\begin{equation}
 (A,B,C)=a^2+xb^2,
 \quad
 \begin{cases}
 a=(xu(x+y),\ 0,\ 1+x^2u),\\
 b=(1+u(x+y+x^2),\ 1+uy,\ uy).
 \end{cases}                                                   \label{eq:DE-A3-ss-pbasis}
\end{equation}
Here the squares are componentwise.  For polynomials over $\mathbf F_2$,
$a,b$ are precisely the two coordinate rows of the three sections of
$V=F_*\OO_E(2qO)$ in the rational basis $(1,x)$, after the indicated
identification of the Frobenius base with $E$.

Let $f_0=A,f_1=B,f_2=C$ denote those three sections of $V$.  On the
Frobenius base consider
\[
\begin{aligned}
 \Psi:V^3&\longrightarrow(\det V)^3,\\
 (v_0,v_1,v_2)&\longmapsto
 \begin{pmatrix}
 f_0\wedge v_1+f_1\wedge v_0\\
 f_0\wedge v_2+f_2\wedge v_0\\
 f_1\wedge v_2+f_2\wedge v_1
 \end{pmatrix}.
\end{aligned}
\]
The second-factor Cartier products in the basis $(1,z,w)$ are
$C(z\,dz)=dz$, $C(w\,dz)=z\,dz$, and $C(zw\,dz)=w\,dz$;
the three square products have Cartier zero.  Consequently $H^0(\Psi)$
is exactly $C_{s_q}$, in the tensor basis with indices $3i+j$.
This identifies the sheaf calculation with the required operator,
including its ordering.

We determine the rank-drop scheme of the three sections of $V$.
Put $\beta=1+uy$.  If $\beta\ne0$, the two rows in
\eqref{eq:DE-A3-ss-pbasis} have rank one precisely when
\[
                 xu(x+y)=0,\qquad 1+x^2u=0.
\]
The second equation forces $x\ne0$.  The first then gives $y=x$;
the elliptic equation gives $x^2+x+1=0$.  Conversely, because
$q/2$ is divisible by three, both points
\[
          Q_\zeta=(\zeta,\zeta),\qquad
          Q_{\zeta^2}=(\zeta^2,\zeta^2),\qquad \zeta^2+\zeta+1=0,
\]
satisfy these equations and have $\beta\ne0$.
If $\beta=0$, then $uy=1$, and the remaining minor is
\[
 a_0b_2+a_2b_0=1+u.
\]
Indeed $1+u=x^3u^2$ follows from the elliptic equation and $uy=1$.
Vanishing of this minor would force $u=1$ and $x=0$, which contradicts
$u=x^{q/2-2}$ with $q\ge6$.  Thus there are no other affine rank drops.
At infinity the leading poles of $A$ and $B$ have different parities,
$2q-1$ and $2q$, and their leading coefficients are one.  They give
an invertible two-by-two matrix in the local Frobenius basis
$(t^{-2q},t^{-2q+1})$.  There is no rank drop there.
At either $Q_\lambda$, $\lambda\in\{\zeta,\zeta^2\}$, the factor
$x+y$ has a simple zero: $\partial(x+y)=1+x^2=\lambda\ne0$.
It follows that the rank-drop ideal is the maximal ideal at each of
these two points, with length one.

Write $K=\ker\Psi$.  Away from these two points, the map
\[
             \End^0(V)\longrightarrow K,\qquad
                    T\longmapsto(Tf_0,Tf_1,Tf_2)
\]
is an isomorphism; here $\End^0$ means trace zero.  At a rank-drop
point choose a uniformizer $t$ and bases so that
$(f_0,f_1,f_2)=(t e_1,e_2,0)$.
Writing $v_0=(a,b)$, $v_1=(c,d)$, $v_2=(e,f)$ gives
\[
                    \Psi(v_0,v_1,v_2)=(a+td,tf,e).
\]
Its kernel consists of $((td,b),(c,d),0)$.
The trace-zero endomorphisms give the submodule in which $b$ is
divisible by $t$.  Thus $K/\End^0(V)$ has length one at each point,
and $\operatorname{coker}\Psi$ also has length one at each point.
In particular $\deg K=2$.

The perfect trace pairing identifies
$(\End^0 V)^\vee=\End(V)/\OO_E\operatorname{id}$.
The local calculation shows that a class in this dual lies in $K^\vee$
exactly when it annihilates the nilpotent endomorphism mapping the
missing direction $e_1$ to the image line $k e_2$.  Equivalently its
representative preserves that image line.  At $Q_\lambda$ the coordinate
row $a$ vanishes, so the image line in the rational basis $(1,x)$ is
$k x$.  By Lemma~\ref{lem:DE-A3-ss-Frobenius-endomorphisms}, a global
class is $\alpha[\operatorname{mult}_x]+\beta[\partial]$, with
$\alpha,\beta\in k$.  The two matrices displayed there show that
preserving this line imposes
\[
                         \alpha\lambda+\beta=0.
\]
The conditions for $\lambda=\zeta$ and $\zeta^2$ are independent.
Hence $H^0(E,K^\vee)=0$.  Serre duality and Riemann--Roch give
$h^1(K)=0$ and $h^0(K)=\deg K=2$.
The two bundles $V$ and $\det V$ have no $H^1$, so the exact sequences
for $\Psi$ identify the cokernel of $H^0(\Psi)$ with the two-dimensional
space of values of its torsion cokernel.  Its rank is therefore $6q-2$.

It remains to check stable rank, rather than merely first rank.
At $Q_\lambda$ the row $b$ is $(0,\lambda,\lambda^2)$, so the local
cokernel functional on an output triple $(U,V,W)$ is
\[
              \ell_\lambda(U,V,W)
                   =\lambda^2 U(Q_\lambda)+\lambda V(Q_\lambda).
\]
The global kernel contains $s_q$ and $\partial_1s_q$, by the Cartier
identities of Lemma~\ref{lem:ss-rectangular-even-bound}.
Their triples of values at $Q_\lambda$ are
\[
              (A,B,C)=(0,1,\lambda^2),\qquad
              (\partial A,\partial B,\partial C)
                        =(0,\lambda^2,\lambda).
\]
Thus the map from the kernel to the cokernel, in these bases, is
\[
 \begin{pmatrix}
    \ell_\zeta(s_q)&\ell_\zeta(\partial_1s_q)\\
    \ell_{\zeta^2}(s_q)&\ell_{\zeta^2}(\partial_1s_q)
 \end{pmatrix}
       =\begin{pmatrix}\zeta&1\\\zeta^2&1\end{pmatrix},
 \qquad \det=1.
\]
In particular the two kernel vectors are independent and their span
meets the image of $C_{s_q}$ trivially.  Since all matrix coefficients
are in $\mathbf F_2$, the Frobenius twists in the second iterate leave
this matrix unchanged.  Therefore its image is stable and
$\rank C_{s_q}^2=\rank C_{s_q}=6q-2$.
The nonzero iterated-rank minors give a nonempty rank open in $|N_q|$.
The first-degree and second-degree line bundles are very ample, so
Bertini gives a nonempty smooth open.  The two opens intersect.
Lemma~\ref{lem:ss-rectangular-even-bound} supplies the matching upper
bound and the identification with relative stable rank.
\end{proof}

\begin{theorem}[All rectangular product polarizations, without a valuation restriction]
\label{thm:DE-A3-all-ss-rectangles}
\status{Proved; root, reduction-descent, and monodromy full independent proof and certificate audits passed.}
Let $E:y^2+y=x^3$ over a perfect field of characteristic two, let
$a,b$ be arbitrary positive integers, and put
\[
 A=E\times E,\qquad M=\OO_E(aO)\boxtimes\OO_E(bO),\qquad L=M^3.
\]
Define $\epsilon(c)=1$ when $c$ is even and $\epsilon(c)=0$ when it
is odd.  For every $d\ge1$ the geometric generic relative stable rank
of a smooth member of $|L^d|$ is
\[
                  r_d=9abd^2-1-\epsilon(ad)-\epsilon(bd).
\]
Its relative coherent dimension is $9abd^2-1$ and its genus is
$9abd^2+1$.  For the embedding by $L$, the same formula applies to
ambient degree-$d$ hypersurfaces whenever
$d\ge\max\{1,\operatorname{reg}(\mathcal I_A)-1\}$.
\end{theorem}
\begin{proof}
Work over an algebraic closure.  If both $ad,bd$ are odd, the separable
polarization theorem, Theorem~\ref{thm:DE-A3-separable-ranks}, gives
$9abd^2-1$.  If both are even, the already proved transverse union
argument in Theorem~\ref{thm:ss-rectangular-all-degrees} gives
$9abd^2-3$.  Thus suppose exactly one is even, say $ad$.
If $v_2(ad)=1$, Lemma~\ref{lem:ss-rectangular-base-norm} gives
$9abd^2-2$.
If $4\mid ad$, put $q=3ad/2$.  It is divisible by six.
Lemma~\ref{lem:DE-A3-ss-uniform-sparse} gives a sharp smooth curve
in $|\OO_E(3adO)\boxtimes\OO_E(3O)|$.
Choose an odd-degree \'etale endomorphism of the first factor of degree
$bd$, and the identity on the second.  Its norm has numerical class
$(3ad,3bd)$.  The complete pair-jet and triple-value incidence proof
of Lemma~\ref{lem:ss-rectangular-base-norm} applies unchanged: the first
factor now has degree at least twelve, and the second still has degree
three.  It uses only separation in degrees at least six and three.
A general sharp source curve therefore has integral nodal image.
Its genus increases from $9ad+1$ to $9abd^2+1$ by exactly
$9ad(bd-1)$ nodes, so its stable rank increases by the same number,
from $9ad-2$ to $9abd^2-2$.
A target translation puts the image in the exact system, because the
polarization of the target ample bundle is surjective on geometric
points.  The smooth rank open then gives the stated lower bound.
The upper bound is Lemma~\ref{lem:ss-rectangular-even-bound}.
This proves every parity case and every valuation.
The dimension, genus, and ambient restriction assertions follow from
the same Riemann--Roch, adjunction, and regularity arguments already
proved in Theorem~\ref{thm:ss-rectangular-all-degrees}.
\end{proof}

\begin{remark}[Scope and retained verification]
\label{rem:DE-A3-all-ss-rectangles-checks}
\status{Proved; root, reduction-descent, and monodromy full independent proof and certificate audits passed.}
The uniform sheaf argument, including the two local modifications and
the invertible kernel--cokernel pairing, is the proof of the infinite
family.  The preceding finite sparse-template searches are retained
as discovery records, including every rejected template and every
nonsharp degree.  In particular template $62556$ fails at many
integers not divisible by six; no extrapolation from its successful
doubling sequence is used.  The exact verification program
\path{agents/obstructions_crystals/code/ss_uniform_rectangular_checks.py}
and its adjacent output check the polynomial identities, the two
geometric points, and eight independently reconstructed finite matrices.
The author run passed in $0.1617353309993632$ seconds; the root reviewer
read the entire program and reran it, preserving the author output,
in $0.154388964001555$ seconds.  A second independent full proof and
program audit passed, with an isolated rerun in
$0.1660513059978257$ seconds.  This result closes the product-polarization valuation
restriction, while making no assertion for arbitrary inseparable
polarizations on a nonproduct supersingular surface or for arbitrary
smooth projective varieties.
\end{remark}
